El uso masivo de herramientas como ChatGPT o GitHub Copilot ha hecho que conseguir la solución a cualquier problema de programación sea instantáneo. Cuando estas tecnologías se usan para aprender sin ningún tipo de límite pedagógico, aparece lo que llamamos "facilismo algorítmico". Los estudiantes suelen dejar que la IA resuelva todo el problema y se limitan a copiar y pegar el código en sus proyectos \cite{chen2024plagiarism}. Este atajo evita el esfuerzo mental que es realmente necesario para entender cómo programar. Estudios recientes muestran que esta dependencia no solo frena el pensamiento crítico, sino que crea una "ilusión de competencia": el alumno tiene un código que funciona, pero no entiende por qué \cite{prather2024widening}. A la larga, esto se traduce en un peor rendimiento y en una pérdida de aprendizaje que se hace evidente al llegar a cursos más difíciles.